% !TEX encoding = UTF-8 Unicode

%%%%%%%%%%%%%%%%%%%%%%%%%%%%%%%%%%%%%%%%%%%%%%%%%%%%%%%%%%%%%%%
%                                                             %
% MINIATURAS MATEM\'ATICAS DE LA GACETA DE LA RSME            %
%                                                             %
%%%%%%%%%%%%%%%%%%%%%%%%%%%%%%%%%%%%%%%%%%%%%%%%%%%%%%%%%%%%%%%

\documentclass[autocontact]{gaceta}

%---------------------
% La Gaceta se har\'a con esto activado, 
% pero no hace falta ahora si no quieres:
\usepackage[T1]{fontenc}
\usepackage{lmodern}

% Puedes usar esto (u otra codificaci\'on) si quieres:
\usepackage[utf8]{inputenc}
\usepackage[spanish]{babel}
%\decimalpoint
%---------------------

%---------------------
% Esto ya lo carga el estilo de La Gaceta:
%\usepackage{amsmath, amsthm, amssymb}
%\usepackage{url} % <-- Para p\'aginas web o similar: \url{...}
%\usepackage{graphicx}
%---------------------

%---------------------
% Carga esto para gr\'aficos si lo necesitas:
% \usepackage{wrapfig}
%---------------------

%---------------------
% Carga esto para gr\'aficos si lo necesitas:
%\usepackage{tikz}
%\usepackage{pgfplots}
%---------------------

%---------------------
% <<< ESTO SE AJUSTAR\'A AL EDITAR CADA N\'UMERO DE LA GACETA >>>
%\provisional
\newcounter{pageforminiatura}
\setcounter{pageforminiatura}{144} % Siempre ser\'a p\'agina par
\setcounter{page}{1} % No se usa (dejarlo como 1)
\journame{La Gaceta de la RSME}
\yearofpublication{2019}
\volume{22}
\issuenumber{1}
%---------------------

%---------------------
% Deja esto as\'{\i}:
%\belongstopart{Art\'{\i}culos} 
%\belongstopart{Secciones} 
%\belongstopart{Noticias de la Sociedad} 
%\belongstopart{Actualidad} 
\nopart
%---------------------

%---------------------
% Usa este tipo de definiciones si las necesitas:
%\theoremstyle{plain}
%\newtheorem*{prop}{Proposici\'on}
%\newtheorem*{lem}{Lema}
% \newtheorem*{teo}{Teorema}
%---------------------


%---------------------
% Deja esto as\'{\i}: 
\supertitle{\large Miniaturas matem\'aticas de \thejourname\\[1mm]
\hfill\footnotesize
Vol.~\thevolume\ (\theyearofpublication), 
N\'um.~\theissuenumber, P\'ag.~\thepageforminiatura\theprovisional}
%
\noeditor
%---------------------

%---------------------
\title{Cálculo del determinante de Vandermonde\\con ayuda de un polinomio}
%\author{} % Poner \author{} o no poner nada, da igual.
% Para ahorrar espacio, el autor solo se pone en la firma:
\contact{Juan Luis Varona, Dpto. de Matemáticas y Computación,
Universidad de La Rioja}{jvarona@unirioja.es}
{http://www.unirioja.es/cu/jvarona/}
%---------------------

%%%%%%%%%%%%%%%%%%%%%%%%%%%%%%%%%%%%%%%%%%%%%%%%%%%%%%%%%%%%%%%%%%%%%%%
\begin{document}
%%%%%%%%%%%%%%%%%%%%%%%%%%%%%%%%%%%%%%%%%%%%%%%%%%%%%%%%%%%%%%%%%%%%%%%

%%%%%%%%%%%%%%%%%%%%%%%%%%%%%%%%%%%%%%%%%%%%%%%%%%%%%%%%%%%%%%%%%%%%%%%
% Deja esto as\'{\i}: 
\previousvspace{-1.8cm}
\makesupertitle
\previousvspace{-0.25cm}
\maketitle
\vspace*{-0.4cm}
\thispagestyle{empty} % <<<<< Importante
%%%%%%%%%%%%%%%%%%%%%%%%%%%%%%%%%%%%%%%%%%%%%%%%%%%%%%%%%%%%%%%%%%%%%%%

%% A PARTIR DE AQU\'I HAY LIBERTAD PARA EL AUTOR:

Se denomina determinante de Vandermonde de orden $n$ a
\[
  V_n = 
  \begin{vmatrix}
  1 & a_{1} & a_{1}^{2} & \ldots & a_{1}^{n-1} \\ 
  1 & a_{2} & a_{2}^{2} & \ldots & a_{2}^{n-1} \\ 
  1 & a_{3} & a_{3}^{2} & \ldots & a_{3}^{n-1} \\ 
  \vdots & \vdots & \vdots & \ddots & \vdots \\ 
  1 & a_{n} & a_{n}^{2} & \ldots & a_{n}^{n-1} 
  \end{vmatrix}.
\]
Es muy conocido que el valor de este determinante es
\begin{equation}
\label{eq:Vn}
  V_n = \prod_{1\leq i < j\leq n}(a_{j}-a_{i}).
\end{equation}
Las demostraciones habituales de~\eqref{eq:Vn} pasan por precisas pero sencillas manipulaciones de $V_n$ aplicando algunas propiedades típicas de los determinantes (por ejemplo, que si a una columna se le resta otra multiplicada por una constante, el valor del determinante no cambia).
Damos aquí una demostración alternativa que usa propiedades elementales de los polinomios y que, pese a ser rápida y elegante, es muy poco conocida (aunque, posiblemente, numerosos matemáticos la han descubierto por su cuenta).
% Posiblemente convendría decir que, sin pérdida de generalidad,
Podemos suponer que los $a_j$ son distintos pues, si no,~\eqref{eq:Vn} es trivial.

Definamos $P_{n-1}(x)$ añadiendo a $V_{n-1}$ una fila y una columna extra:
% como se muestra a continuación:
\[
  P_{n-1}(x) = 
  \begin{vmatrix}
   &  &  &  &  &            a_{1}^{n-1} \\ 
   &  &  &  &  &            a_{2}^{n-1} \\ 
   &  & \hspace{-1ex}V_{n-1}\hspace{-1ex} &  &  &     a_{3}^{n-1} \\ 
   &  &  &  &  &            \vdots \\ 
   &  &  &  &  &            a_{n-1}^{n-1} \\ 
  1 & x & x^{2} & \ldots &  x^{n-2} & x^{n-1} 
  \end{vmatrix}.
\]
%Por la propiedades de los determinantes, es claro que
%\begin{itemize}
%\item $P_{n-1}(x)$ es un polinomio de grado $n-1$ y coeficiente director $V_{n-1}$ (basta desarrollar por la última fila).
%\item $P_{n-1}(x)$ se anula en $x = a_1,a_2,\ldots,a_{n-1}$ (el correspondiente determinante tiene dos filas iguales).
%\end{itemize}
Es claro que $P_{n-1}(x)$ es un polinomio de grado $n-1$ y coeficiente director $V_{n-1}$ (basta desarrollar por la última fila);
además, $P_{n-1}(x)$ se anula en $x = a_1,a_2,\ldots,a_{n-1}$ (el correspondiente determinante tiene dos filas iguales).
Así pues,
\[
  P_{n-1}(x) = V_{n-1} \cdot (x-a_1)(x-a_2) \cdots (x-a_{n-1}).
\]
Pero
\[
  V_n = P_{n-1}(a_n) = V_{n-1} \cdot (a_n-a_1)(a_n-a_2) \cdots (a_n-a_{n-1}),
\]
y esto prueba~\eqref{eq:Vn} por inducción.

\enlargethispage{2mm}

\medskip

%%%%%%%%%%%%%%%%%%%%%%%%%%%%%%%%%%%%%%%%%%%%%%%%%%%%%%%%%%%%%%%%%%%%%%%
\end{document}
%%%%%%%%%%%%%%%%%%%%%%%%%%%%%%%%%%%%%%%%%%%%%%%%%%%%%%%%%%%%%%%%%%%%%%%
